\chapter{Performance}

\section{Metriche di Riferimento}

Le seguenti metriche sono state misurate su un dataset reale dell'RNA composto da circa 165 file XML:

\begin{table}[h]
\centering
\begin{tabular}{@{}ll@{}}
\toprule
\textbf{Metrica} & \textbf{Valore} \\
\midrule
Velocità di parsing & $\sim$50.000 record/sec (8 worker) \\
Utilizzo memoria per worker & $\sim$200 MB \\
Rapporto compressione & $\sim$10x vs XML grezzo \\
\bottomrule
\end{tabular}
\caption{Benchmark su dataset RNA reale}
\end{table}

\section{Scalabilità Orizzontale}

La performance scala linearmente con il numero di worker fino al numero di core fisici disponibili:

\begin{figure}[h]
\centering
\begin{tikzpicture}
    \begin{axis}[
        width=10cm,
        height=6cm,
        xlabel={Numero Worker},
        ylabel={Throughput (record/sec)},
        xmin=1, xmax=8,
        ymin=0, ymax=55000,
        xtick={1,2,4,6,8},
        ytick={0,10000,20000,30000,40000,50000},
        legend pos=north west,
        grid=major,
        grid style={gray!30}
    ]
    \addplot[color=primaryblue, mark=*, thick] coordinates {
        (1, 7000)
        (2, 14000)
        (4, 27000)
        (6, 40000)
        (8, 50000)
    };
    \addlegendentry{Misurato}
    
    \addplot[color=gray, dashed] coordinates {
        (1, 7000)
        (8, 56000)
    };
    \addlegendentry{Lineare ideale}
    \end{axis}
\end{tikzpicture}
\caption{Scalabilità throughput vs numero worker}
\end{figure}

\begin{infobox}[Nota sulle Performance]
Il leggero scostamento dalla linearità ideale è dovuto a:
\begin{itemize}
    \item Overhead di serializzazione/deserializzazione tra processi
    \item Contesa I/O su disco per scritture Parquet concorrenti
    \item Scheduling overhead del sistema operativo
\end{itemize}
\end{infobox}

\section{Compressione Parquet}

Il formato Parquet offre vantaggi significativi rispetto all'XML:

\begin{table}[h]
\centering
\begin{tabular}{@{}lll@{}}
\toprule
\textbf{Aspetto} & \textbf{XML} & \textbf{Parquet} \\
\midrule
Storage size & 1 GB & $\sim$100 MB \\
Tempo lettura completa & $>$60s & $<$5s \\
Query selettive & Full scan & Partition pruning \\
Tipizzazione & Assente & Nativa \\
\bottomrule
\end{tabular}
\caption{Confronto XML vs Parquet}
\end{table}

\section{Ottimizzazioni Chiave}

\subsection{Streaming Parser}

L'utilizzo di \texttt{iterparse()} evita di caricare l'intero DOM in memoria:

\begin{itemize}
    \item Memory footprint costante indipendentemente dalla dimensione file
    \item Possibilità di processare file multi-gigabyte su macchine con RAM limitata
\end{itemize}

\subsection{Batch Writing}

La scrittura in batch da 10.000 record bilancia:
\begin{itemize}
    \item \textbf{Overhead I/O}: meno syscall rispetto a scritture singole
    \item \textbf{Memory usage}: batch sufficientemente piccoli da non saturare RAM
\end{itemize}

\subsection{Aggressive Cleanup}

Il pattern di cleanup post-parsing previene memory leak:

\begin{lstlisting}[language=Python]
elem.clear()
while elem.getprevious() is not None:
    del elem.getparent()[0]
\end{lstlisting}

Senza questo cleanup, \texttt{lxml} mantiene riferimenti all'intero documento ancestry.
